\section{The deterministic M\"obius corner as one Walsh spectrum}
\label{sec:tpc43-Walsh-corner}

The equality alias admits a sharper deterministic description than the
annealed theorem.  For nonzero atoms, let
\(u_{\alpha,j}\) be the squarefree fused label in
\eqref{eq:tpc43-fused-zero-label}.  For two rows at the same orbit time,
define their squarefree symmetric--difference kernel
\begin{equation}
 \kappa_{\alpha,\beta,j}
 :=\frac{u_{\alpha,j}u_{\beta,j}}
 {(u_{\alpha,j},u_{\beta,j})^2}.
 \label{eq:tpc43-pair-kernel}
\end{equation}
It is squarefree and
\begin{equation}
 \chi_{u_{\alpha,j}}(\varepsilon)
 \chi_{u_{\beta,j}}(\varepsilon)
 =\chi_{\kappa_{\alpha,\beta,j}}(\varepsilon).
 \label{eq:tpc43-pair-character}
\end{equation}

For \(\kappa>1\), let \(C(\kappa)\) be the sum of all ordered unequal
row coefficients, on both physical sides, whose pair kernel is
\(\kappa\):
\begin{align}
 C(\kappa)
 :={}&\sum_{\gamma,j}
 \sum_{\substack{\alpha\ne\beta\\
 \kappa_{\alpha,\beta,j}=\kappa}}
 a^L_{\alpha,\gamma,j,0}
 \overline{a^L_{\beta,\gamma,j,0}}
 \mu^2(n_{\alpha,j,0})\mu^2(n_{\beta,j,0})
 \notag\\
 &+\text{the symmetric right contribution}.
 \label{eq:tpc43-Walsh-spectrum-coefficient}
\end{align}
Pairing \((\alpha,\beta)\) with \((\beta,\alpha)\) shows that
\(C(\kappa)\in\R\).

\begin{theorem}[Exact four--to--one M\"obius compression]
\label{thm:tpc43-four-to-one-compression}
The complete equality energy has the Walsh expansion
\begin{equation}
 \boxed{
 Z_0(\varepsilon)
 =\cD_0+
  \sum_{\kappa>1}C(\kappa)\chi_\kappa(\varepsilon).}
 \label{eq:tpc43-exact-Walsh-expansion}
\end{equation}
At the physical all--minus point,
\begin{equation}
 \boxed{
 \|\mathscr T_0^L\|^2+\|\mathscr T_0^R\|^2
 =\cD_0+
  \sum_{\kappa>1}\mu(\kappa)C(\kappa).}
 \label{eq:tpc43-physical-kernel-spectrum}
\end{equation}
Moreover,
\begin{equation}
 \E_\varepsilon
 |Z_0(\varepsilon)-\cD_0|^2
 =\sum_{\kappa>1}|C(\kappa)|^2.
 \label{eq:tpc43-Walsh-Parseval}
\end{equation}
\end{theorem}

\begin{proof}
Expand the two squared Hilbert norms defining \(Z_0\).  By
\eqref{eq:tpc43-pair-character}, every ordered row pair contributes the
Walsh character indexed by its symmetric--difference kernel.  The
constant character occurs exactly when
\(u_{\alpha,j}=u_{\beta,j}\), which by
\cref{lem:tpc43-fused-label-injectivity} means \(\alpha=\beta\).  Its
coefficient is therefore \(\cD_0\), proving
\eqref{eq:tpc43-exact-Walsh-expansion}.

At \(\varepsilon_p=-1\),
\begin{equation}
 \chi_\kappa(-\one)=(-1)^{\omega(\kappa)}=\mu(\kappa).
 \label{eq:tpc43-character-at-minus}
\end{equation}
For a contributing pair, this also gives the literal identity
\begin{align}
 &\mu(d_\alpha)\mu(n_{\alpha,j,0})
  \mu(d_\beta)\mu(n_{\beta,j,0})
 \notag\\
 &\qquad
 =\mu(u_{\alpha,j})\mu(u_{\beta,j})
 =\mu(\kappa_{\alpha,\beta,j}).
 \label{eq:tpc43-literal-four-to-one}
\end{align}
Thus \eqref{eq:tpc43-physical-kernel-spectrum} is precisely the physical
four--M\"obius expansion, not a surrogate.  Finally, Walsh orthogonality
and the absence of a nonzero constant term give
\eqref{eq:tpc43-Walsh-Parseval}.
\end{proof}

The theorem replaces a row--pair family of four M\"obius signs by one
signed spectral sum over squarefree kernels.  The coefficient
\(C(\kappa)\) still contains all masks and all representations of the
same kernel; no cancellation estimate for that sum is asserted.

There is an exact small--prime filtration.  Fix \(z\ge2\), average the
signs \(\varepsilon_p\) for \(p\le z\), and set
\(\varepsilon_p=-1\) for \(p>z\).  Let \(P^-(n)\) denote the least prime
factor of \(n>1\).

\begin{corollary}[Rough--kernel filtration]
\label{cor:tpc43-rough-kernel-filtration}
One has
\begin{equation}
 \boxed{
 \E_{\varepsilon_p:\,p\le z}
 Z_0(\varepsilon_{\le z},-\one_{>z})
 =\cD_0+
  \sum_{\substack{\kappa>1\\P^-(\kappa)>z}}
  \mu(\kappa)C(\kappa).}
 \label{eq:tpc43-rough-filtration}
\end{equation}
\end{corollary}

\begin{proof}
The average of \(\chi_\kappa\) is zero if any prime
\(p\le z\) divides \(\kappa\).  Otherwise every prime in \(\kappa\)
keeps the value \(-1\), and the character equals \(\mu(\kappa)\).
Insert this into \eqref{eq:tpc43-exact-Walsh-expansion}.
\end{proof}

Before any relevant prime is averaged one has the physical corner;
after every relevant prime has been averaged, only \(\cD_0\) remains.
A future
de--randomization may therefore attack the successive prime increments
or the final rough spectrum.  The identity itself does not bound either.
